
\documentclass{article}
\usepackage{colortbl}
\usepackage{makecell}
\usepackage{multirow}
\usepackage{supertabular}

\begin{document}

\newcounter{utterance}

\centering \large Interaction Transcript for game `cladder', experiment `full\_v1.5\_default', episode 8201 with reflexion{-}Qwen3.5{-}9B.
\vspace{24pt}

{ \footnotesize  \setcounter{utterance}{1}
\setlength{\tabcolsep}{0pt}
\begin{supertabular}{c@{$\;$}|p{.15\linewidth}@{}p{.15\linewidth}p{.15\linewidth}p{.15\linewidth}p{.15\linewidth}p{.15\linewidth}}
    \# & \multicolumn{2}{c}{Player} && \multicolumn{2}{c}{Game Master} \\
    \hline

    \theutterance \stepcounter{utterance}  
    & & & \multicolumn{4}{p{0.6\linewidth}}{
        \cellcolor[rgb]{0.9,0.9,0.9}{
            \makecell[{{p{\linewidth}}}]{
                \texttt{\tiny{[P1$\langle$GM]}}
                \texttt{You have attempted this same task before and did not succeed. These are the plans you wrote after those attempts.} \\
\texttt{Trial \#1: New plan: I failed because I attempted to calculate the counterfactual probability \$P(gyzp \textbackslash{}mid do(xevo=0))\$ by reasoning about the causal graph structure and the intervention's effect on downstream variables, but I did not actually perform the necessary calculation using the provided conditional probabilities. Specifically, I stopped mid{-}sentence while trying to determine the new distribution of \$tijw\$ under the intervention \$do(xevo=0)\$. Since \$xevo\$ is a child of \$Tijv\$, setting \$xevo=0\$ does not change the probability of \$tijw\$ (which is determined by \$Tijv\$); thus, \$P(tijw=1 \textbackslash{}mid do(xevo=0))\$ remains \$0.28\$. I should have immediately calculated \$P(gyzp \textbackslash{}mid do(xevo=0))\$ as \$(0.28 \textbackslash{}times 0.14) + (0.72 \textbackslash{}times 0.76) = 0.0392 + 0.5472 = 0.5864\$. Comparing this to the actual \$P(gyzp \textbackslash{}mid xevo=1) = 0.1228\$, the probability is significantly lower if \$xevo\$ is not present. Next time, I will explicitly calculate both probabilities using the law of total probability with the correct marginal distributions before drawing a conclusion.} \\
\\ 
\\ 
\texttt{You are an expert in causal inference. The following question is not a typical commonsense query, but rather a meticulously designed question created by a professor specializing in causal inference, intended to assess the students' mastery of the course content.} \\
\texttt{Imagine a self{-}contained, hypothetical world with only the following conditions, and without any unmentioned factors or causal relationships: Tijv has a direct effect on tijw and xevo. Xevo has a direct effect on gyzp. Tijw has a direct effect on gyzp. For those who are not xevo and are not tijw, the probability of gyzp is 76\%. For those who are not xevo and are tijw, the probability of gyzp is 76\%. For those who are xevo and are not tijw, the probability of gyzp is 12\%. For those who are xevo and are tijw, the probability of gyzp is 14\%. For those who are not xevo, the probability of tijw is 28\%. For those who are xevo, the probability of tijw is 14\%. For those who are xevo, would it be more likely to see gyzp if the individual was not xevo?} \\
\texttt{Start your answer with 'Yes' or 'No', followed by additional reasoning or evidence to support your explanation.} \\
            }
        }
    }
     \\ \\

    \theutterance \stepcounter{utterance}  
    & \multicolumn{4}{p{0.6\linewidth}}{
        \cellcolor[rgb]{0.561,0.737,0.561}{
            \makecell[{{p{\linewidth}}}]{
                \texttt{\tiny{[P1$\rangle$GM]}}
                \texttt{No} \\
\\ 
\texttt{**Reasoning:**} \\
\\ 
\texttt{To determine if \$gyzp\$ is more likely under the intervention \$do(xevo=0)\$ compared to the observed state \$xevo=1\$, we must calculate the probability of \$gyzp\$ in both scenarios using the provided conditional probabilities and the causal structure.} \\
\\ 
\texttt{**1. Define the Causal Structure and Probabilities:**} \\
\texttt{Based on the prompt:} \\
\texttt{*   **Parents of \$gyzp\$:** \$xevo\$ and \$tijw\$.} \\
\texttt{*   **Parents of \$tijw\$:** \$xevo\$ (implied by the conditional probabilities given \$xevo\$).} \\
\texttt{*   **Parents of \$xevo\$:** \$Tijv\$ (implied by the prompt stating "Tijv has a direct effect on... xevo").} \\
\\ 
\texttt{The conditional probabilities for \$gyzp\$ are:} \\
\texttt{*   \$P(gyzp=1 \textbackslash{}mid xevo=0, tijw=0) = 0.76\$} \\
\texttt{*   \$P(gyzp=1 \textbackslash{}mid xevo=0, tijw=1) = 0.76\$} \\
\texttt{*   \$P(gyzp=1 \textbackslash{}mid xevo=1, tijw=0) = 0.12\$} \\
\texttt{*   \$P(gyzp=1 \textbackslash{}mid xevo=1, tijw=1) = 0.14\$} \\
\\ 
\texttt{The conditional probabilities for \$tijw\$ are:} \\
\texttt{*   \$P(tijw=1 \textbackslash{}mid xevo=0) = 0.28\$} \\
\texttt{*   \$P(tijw=1 \textbackslash{}mid xevo=1) = 0.14\$} \\
\\ 
\texttt{**2. Calculate \$P(gyzp \textbackslash{}mid xevo=1)\$ (Observed Probability):**} \\
\texttt{We use the Law of Total Probability, conditioning on \$tijw\$:} \\
\texttt{\$\$P(gyzp \textbackslash{}mid xevo=1) = P(gyzp \textbackslash{}mid xevo=1, tijw=0)P(tijw=0 \textbackslash{}mid xevo=1) + P(gyzp \textbackslash{}mid xevo=1, tijw=1)P(tijw=1 \textbackslash{}mid xevo=1)\$\$} \\
\\ 
\texttt{*   \$P(tijw=0 \textbackslash{}mid xevo=1) = 1 {-} 0.14 = 0.86\$} \\
\texttt{*   \$P(tijw=1 \textbackslash{}mid xevo=1) = 0.14\$} \\
\\ 
\texttt{\$\$P(gyzp \textbackslash{}mid xevo=1) = (0.12 \textbackslash{}times 0.86) + (0.14 \textbackslash{}times 0.14)\$\$} \\
\texttt{\$\$P(gyzp \textbackslash{}mid xevo=1) = 0.1032 + 0.0196 = 0.1228\$\$} \\
\\ 
\texttt{**} \\
            }
        }
    }
    & & \\ \\

    \theutterance \stepcounter{utterance}  
    & & & \multicolumn{2}{p{0.3\linewidth}}{
        \cellcolor[rgb]{0.9,0.9,0.9}{
            \makecell[{{p{\linewidth}}}]{
                \texttt{\tiny{[GM$|$GM]}}
                \texttt{No} \\
            }
        }
    }
    & & \\ \\

    \theutterance \stepcounter{utterance}  
    & & & \multicolumn{2}{p{0.3\linewidth}}{
        \cellcolor[rgb]{0.9,0.9,0.9}{
            \makecell[{{p{\linewidth}}}]{
                \texttt{\tiny{[GM$|$GM]}}
                \texttt{yes} \\
            }
        }
    }
    & & \\ \\

    \theutterance \stepcounter{utterance}  
    & & & \multicolumn{2}{p{0.3\linewidth}}{
        \cellcolor[rgb]{0.9,0.9,0.9}{
            \makecell[{{p{\linewidth}}}]{
                \texttt{\tiny{[GM$|$GM]}}
                \texttt{game\_result = LOSE} \\
            }
        }
    }
    & & \\ \\

    \theutterance \stepcounter{utterance}  
    & & \multicolumn{4}{p{0.6\linewidth}}{
        \cellcolor[rgb]{0.95,0.95,0.95}{
            \makecell[{{p{\linewidth}}}]{
                \texttt{\tiny{[P1$\leftrightarrow$P]}}
                \texttt{Post{-}episode Reflexion (not sent to GM):} \\
\texttt{New plan: I failed because I correctly identified that the intervention \$do(xevo=0)\$ severs the causal link from \$xevo\$ to \$tijw\$, leaving \$P(tijw)\$ at its baseline \$0.28\$, but I stopped my reasoning before completing the final calculation for \$P(gyzp \textbackslash{}mid do(xevo=0))\$. Specifically, I did not sum the weighted probabilities for the two possible states of \$tijw\$ under the intervention: \$(0.28 \textbackslash{}times 0.14) + (0.72 \textbackslash{}times 0.76)\$. This calculation yields \$0.0392 + 0.5472 = 0.5864\$. Since \$0.5864\$ is significantly greater than the observed probability \$0.1228\$, the correct answer is "Yes". Next time, I will explicitly perform the arithmetic for the counterfactual probability immediately after establishing the correct marginal distribution for \$tijw\$ to ensure the final comparison is accurate.} \\
            }
        }
    }
    & \\ \\

\end{supertabular}
}

\end{document}
